\section{Introduction}

Modern language-model systems have several ways to use memory. They can answer directly from parameters, retrieve fresh evidence, cache useful facts in external memory, adapt temporarily to a local update, or commit a stable update more durably. The core question in this project is not whether any one of these operations is useful in isolation. It is when a controller should choose one operation instead of another.

This framing matters because different regimes reward different choices. On recurring factual queries, repeated retrieval may be wasteful once a reusable answer has been observed and stored. On updated knowledge, cached memory may become stale, retrieval may lag, and temporary or durable adaptation can become preferable. A controller that sees all of these options under one objective is therefore more expressive than a retrieval gate or a standalone memory writer.

This paper makes two narrow, artifact-backed claims. First, on PopQA, the controller matches retrieval-level quality overall after a conservative non-recurring guardrail fix while reducing repeated retrieval on recurring queries. Second, on a deterministic bundled freshness/update benchmark, the controller improves updated-knowledge behavior over retrieval-only and memory-only baselines while exercising temporary adaptation, rare guarded consolidation, and rollback.

The contribution is therefore a unified multi-timescale routing perspective backed by two complementary experiments:
\begin{itemize}[leftmargin=1.5em]
\item \textbf{PopQA} validates reusable external memory and delayed-utility routing.
\item \textbf{Freshness} validates temporary adaptation, guarded durable consolidation, and rollback under updated knowledge.
\end{itemize}
